\hypertarget{Vedlegg A}{}
\section*{Vedlegg A}
\addcontentsline{toc}{section}{Vedlegg A}
\appendix

Dette vedlegget inneholder implementasjonen av de egenlagde
funksjonene som brukes i Oppgave 3.

% --------------------------------------------------
\subsection*{mySquareV1}

\begin{verbatim}
function [t, out] = mySquareV1(t, f0, d_c, ...
                               bds=[-1,1], pds=0, cent=false);
    #{
        Denne funksjonen genererer firkantpulsen ved å uttrykke
        signalets fase i radianer og bruke modulo 2*pi.

        t: tidsvektor
        f0: fundamental frekvens
        d_c: duty cycle
        bds:  grenser -> bds(1) = lav verdi, bds(2) = høy verdi.
        pds: antall perioder man ønsker: pds <= antall perioder det er i t
        cent: få en sentrert positiv puls omkring origo

        return verdier
        t:  pds=0 -> t_ut = t_in ellers er den tilpaset. [ms]
        out: resulterde verdier for y aksen (amplitude vektor).
    #}

    T = 1/f0;

    if pds != 0;
        t_stop = T*pds;
        idx = find(t >= t_stop, 1)-1;
        t = t(1:idx);
    end;

    d_c = d_c/100; # konverter til desimal
    rads = t*2*pi*f0; # konverter til radianer

    if cent;
        rads = t*2*pi*f0 + (2*pi*d_c)/2;
    end;

    pulse_end_rad = 2*pi*d_c; # når pulsen skal gå lavt igjen
    n_rad = mod(rads, 2*pi); # normaliser radianer til 0-2pi

    out = n_rad; 
    out(n_rad >= pulse_end_rad) = bds(1);
    out(n_rad < pulse_end_rad) = bds(2);

    t = t*1e3; # skaler til ms
end;
\end{verbatim}


% --------------------------------------------------
\subsection*{mySquareV2}

\begin{verbatim}
function [t, out] = mySquareV2(f0, d_c, spp, ...
                               bds=[-1,1], pds=5, sh=0);
    #{
        Denne funksjonen genererer firkantpulsen direkte i
        tidsdomenet ved å bruke modulo av periodetiden T.

        f0: fundamental frekvens
        d_c: duty cycle
        spp: samples per periode
        bds: grenser -> bds(1) = lav verdi, bds(2) = høy verdi.
        pds: antall perioder man ønsker
        shfit: sh i tid gitt i prosent -> t + T*(sh/100)

        return verdier
        t: resulterdne tidsavektor skalert til ms,
        out: resulterde verdier for y aksen (amplitude vektor).
    #}

    d_c = d_c/100; # konverter til desimal
    T = 1/f0; 
    dt = T/spp;
    pulse_end_time = T*d_c; # når pulsen skal gå lavt igjen i tid 

    t = 0:dt:(T*pds-dt);
    t_mod = mod(t + T*(sh/100), T); # normaliser (t + t_shift) til 0-T 

    out = t_mod;
    out(t_mod < pulse_end_time) = bds(2);
    out(t_mod >= pulse_end_time) = bds(1);

    t = t*1e3; # skaler til ms
end;
\end{verbatim}

% --------------------------------------------------
\subsection*{addZerosOrOnes}

\begin{verbatim}
function y = addZerosOrOnes(x, num);
    if x == 0;
        y = zeros(1, num);
    else
        y = ones(1, num);
    end;
end;
\end{verbatim}


% --------------------------------------------------
\subsection*{bitWaveGen}

\begin{verbatim}
function [t, out] = bitWaveGen(bits, T, spp, sh=0)
    #{
        bits: vektor med bits (1 og 0) som skal generere en puls
        T: bit tid (periode) i sekunder
        spp: samples per periode
        shift: sh i tid gitt i prosent -> t + T*(sh/100)

        return verdier
        t: resulterende tidsvektor skalert til ms,
        out: resulterende verdier for y aksen (amplitude vektor).
    #}

    dt = T/spp;
    pds = length(bits); # antall perioder
    sh_num = round(spp*(sh/100)); # antall som skal shiftes

    t = 0:dt:(T*pds-dt);
    out = [];

    first = bits(1); # bit som skal shiftes først

    for i=1:pds;

        if i == 1;
            out = [out, ...
                   addZerosOrOnes(bits(i), spp-sh_num)];
        else
            out = [out, ...
                   addZerosOrOnes(bits(i), spp)];
        end;
    end;

    # legg til shift rest på første bit tilbake i vektoren
    out = [out, addZerosOrOnes(first, sh_num)]; 

    t = t*1e3; # skaler til ms
end;
\end{verbatim}